\subsection{The ordinary ranks of the retained digit classes}
\label{sec:retained-rank}
\hypertarget{source-retained-rank}{}

A successful specialization controls the Gram form on the retained space. Estimate its ordinary dimension before multiplying the segment contributions in \eqref{eq:ordinary-class}.
\begin{established}[Separation and degree reserves on the retained digit class]\label{hyp:retained-digit-degrees}
Work in one compatible digit class. Write \(B=A_0^{\,*}\),
\(M=\lceil m_0\rceil\), \(\delta=4/p^2\) and
\(\epsilon=N_1\mathcal L/k\). The separation of the odd powers
from the transition between the two lists is important here.
We give the numerical details used from the separated regime.
Let \(\Delta<1/6\) be its relative band width parameter, with
\(40\epsilon\le\Delta\).

At the first list
\(p^{d_i}<(1-2\Delta)a/b'\), at the second
\(p^{d_i}>(1+2\Delta)a/b'\). In the enlarged box
\(|b_i/b'-1|\le b_e\le\Delta/2-18\epsilon\),
\((E_0+Q)/a\le\epsilon\) and \(kQ\le\epsilon b'\).

\begin{equation}\label{eq:coarse-split}
 (1+\Delta)A_i<B\quad(i<\sigma),\qquad
 B<(1-\Delta)A_i\quad(i\ge\sigma).
\end{equation}
The prefix sums satisfy

\begin{equation}\label{eq:coarse-prefix}
 \sum_{l<i} A_l\le\delta A_i,\qquad
 \sum_{l<i} C_l\le\delta C_i .
\end{equation}
\end{established}
\begin{proof}
The inequalities for the box are verified on the selected host cycle in the size calculation of Proposition~\ref{prop:host-large-prime-test}. At each interior top the digit and floor identities give \[
 \begin{gathered}
 b_i-kP_i<P_i A_i\le b_i, \\ P_i p^{d_i}=p^G, \\ a-E_0-p^G<Bp^G\le a.
\end{gathered}
\] For the first inequality use
\((1-2\Delta)(1+b_e)(1+\Delta)\le1-\epsilon\).

For the second use
\((1+2\Delta)(1-b_e-\epsilon)(1-\Delta)\ge1\).
Both follow by expanding with the bounds on \(b_e,\Delta\).
For consecutive indices \(P_i/P_{i+1}\ge p^2\), as both
distances are odd. The bounds just used give
the two ratio bounds
\[
       p^2A_i\le2A_{i+1},\qquad p^2C_i\le2C_{i+1},
\]
using \(a_i\in[a,a+2a/(kY)]\) and the digit error \(kP_i\). Summing the resulting geometric progressions gives \eqref{eq:coarse-prefix}.

We may use \(B\ge a/(2Q)\), by the floor bound and
\((E_0+Q)/a\le\epsilon\). With \(4 k^3 QY\le b'\),
\(kY\mathcal L b'\le a\), \(\mathcal L\ge24\), and
\(m_0\le b'/(kQY)\), we have
\[
 kM/B\le4b'/a\le1/30,\qquad M\le\epsilon B.
\]
Also \(k^3\le A_i,C_i,B-M\) by the digit and floor bounds,
and \(k\ge10\) by \(40\epsilon\le\Delta\). These inequalities
apply with absolute size constants, simultaneously on the
enlarged box.
\end{proof}

\begin{proposition}[Ordinary rank of one retained digit class]\label{prop:retained-digit-rank}
Fix a compatible digit class satisfying
Lemma~\ref{hyp:retained-digit-degrees}, using the retained space
in \eqref{eq:ordinary-class}.
Let \(R_{\rm ret}\) be the ordinary rank of the retained space
and \(N_{\rm mon}\) the unrestricted monomial count of this class. Then
\[
 R_{\rm ret}\ge N_{\rm mon}
 (1-e^{-\mathcal L/8})(1-e^{-\mathcal L/3})(1-rM/B).
\]
For \(\mathcal L\ge24\) and \(kM/B\le1/30\), this fraction
is greater than one half. When \(s=0\) there is no loss
from the internal tests.
\end{proposition}

\begin{proof}
\proofstep{The two counts.}
For \(A\in\{B,B-M\}\), put \(V=\mathbb C^k\). The tensor and
its dimension are
\[
 \mathcal T_A=\operatorname{Sym}^{A}V\otimes
       \bigotimes_i(\operatorname{Sym}^{A_i}V\otimes
                         \operatorname{Sym}^{C_i}V^*),\qquad
 T(A)=D_k(A)\prod_i D_k(A_i)D_k(C_i).
\]
In characteristic zero, each Pieri step gives one copy of each
allowed highest weight. For decreasing integer rows write
\(\nu\succ\mu\) when \(\nu_j\ge\mu_j\ge\nu_{j+1}\) wherever
the entries exist. From the zero weight add the initial strip of
size \(A\). At each subsequent pair the allowable transitions are
\[
 \mu\ \longrightarrow\ \nu\ \longrightarrow\ \tau,\qquad
 \nu\succ\mu,\quad |\nu|-|\mu|=A_i,\qquad
 \nu\succ\tau,\quad |\nu|-|\tau|=C_i ,
\]
where \(|\mu|\) denotes the sum of the rows, including negative
entries. These rules follow from their polynomial versions by a
determinant twist and duality. Restriction to \(\mathrm{GL}_r\)
appends one interlacing weight \(\lambda\) below the last weight,
again with multiplicity one. \cite{fulton-harris}

Let \(\mathcal P_A\) denote these entire paths and
\(d_\lambda=\dim V_\lambda\). The terminal dimension and the
independent choice of the last-coordinate exponent \(j_+\)
give the two counts
\begin{equation}\label{eq:retained-path-dimensions}
 \sum_{\pi\in\mathcal P_A}d_{\lambda(\pi)}=T(A),\qquad
 \sum_{j_+=0}^{b_+}D_r(b_+-j_+)=D_k(b_+).
\end{equation}
The second identity partitions degree-\(b_+\) monomials
according to their exponent on the last variable.

\proofstep{The row recursion for a counted path.}
Number rows from zero. Let \(c_i\) be the first row before
the \(i\)-th vector addition, \(c_0=A\), and \(c_s\) the final
one. Let \(w_i\) be the increase in the sum of all other
rows at the addition and \(h_i\) the decrease in the first
row at the removal. Write
\[
 u_i=w_i+h_i,
\]
and let \(R_i\)
be row 1 just after the addition.

Write \(n_i=-\lambda_{r-s+i}\)
and let \(d\) be the sum of \(\lambda\) minus the sum of the
rows below the first row at the end. Deleting the first row
after the initial letter gives a path in dimension \(r\)
with alternating strip sizes \(w_i,C_i-h_i\) and one final
addition of size \(d\) ending at \(\lambda\). Indeed the last
interlacing is exactly such a horizontal addition. Record \((w_i,h_i)\) in the initial list. The tied positions and the endpoint have labels
\[
 (o_i,j_i)=(u_i,h_i),\qquad o_+=d+j_+.
\]
Conversely these data and
the reduced path uniquely recover the large path at a
specified \(A\). Its first rows at the three successive times
of a pair must be
\[
 c_i,\quad c_i+A_i-w_i,\quad c_{i+1}=c_i+A_i-u_i .
\]
Thus sizes determine the restored entries. Their interlacing
with the retained rows is the condition to check for the
converse.

\proofstep{Bounding the upper count by the unrestricted list.}
For the upper bound start with a counted block, level and
allowed path in dimension \(r\). The independent factors in \eqref{eq:retained-path-dimensions}
give the count on this side as the sum of \(d_\lambda D_r(b_+-j_+)\) because each prefix
cutoff takes the sum of the allowed prefix summands in the
successive Pieri decompositions. The triangular lift of levels
in Section~\ref{sec:segment-window} preserves this total rank.
Reconstruct at \(A=B\), where \(c_i=K_i\). The first row of the
reduced path just after its vector must be at most both
\(K_i,K_{i+1}\). At a tie its cutoff gives the former and
\(K_{i+1}\ge K_i\).

In the initial list this first row is at
most
\[
 \sum_{l\le i}w_l\le(1+\delta)A_i,
\]
since a covector
step never increases it. Up through this pair every
\(K_l\ge B-2s\mathcal L B/k\), by \(w_l\le A_l\) and the
bound on \(h_l\). By \eqref{eq:coarse-split} and
\[
 \delta+3\epsilon\le\Delta/2,
\]
both required inequalities
follow. The final cutoff gives \(\lambda_0\le K_s\). The strip
degrees give \(0\le h_i\le C_i\), \(0\le w_i\le A_i\)
also at the ties. These conditions suffice for every inserted
horizontal strip and the final restriction. We have a
weight-preserving injection into the unrestricted count,
proving the upper bound of \eqref{eq:ordinary-class}.

\proofstep{A lower bound by deletion.}
For the lower bound start with \(A=B-M\), so the capacities
reconstructed from the labels satisfy \(K_i=c_i+M\). Retain
only paths satisfying, at every \(i\),
\begin{equation}\label{eq:kept-tail}
 u_i-R_i+\delta(A_i+r)\le \frac{\mathcal L}{k}\max(c_i,A_i),
 \qquad n_i\ge (1-\mathcal L/k)C_i ,
\end{equation}
and \(j_+\le\mathcal L b_+/k\).
We verify that these paths contribute to the block count at
margin \(m_0\). Interlacing and \eqref{eq:coarse-prefix} give
\[
 0\le R_i\le c_i,\qquad
 R_i\le w_i+\delta A_i,\qquad
 0\le n_i\le(1+\delta)C_i .
\]

\proofstep{Keeping enough rows after deletion.}
For the second bound the previous row 1 was at most the sum of
previous vector additions below the first row. For the third,
follow the bottom row just after pair \(i\). Until then its
negative part was at most \(\sum_{l\le i} C_l\). A later vector
step increases each row, a covector step decreases a row at
most to the former next row, and \(\lambda_j\ge\mu_{j+1}\).
Following these comparisons back proves the upper bound. A horizontal strip creates at most one new row of the corresponding sign. After \(i\) pairs, starting with the initial vector, the row counts thus satisfy 
\[
 \#\{a:\lambda_a>0\}\le i+1,\qquad\#\{a:\lambda_a<0\}\le i,
\]
 for the intermediate highest weight \(\lambda\). The inequalities about signs follow
from the same interlacing, in particular row 1 is nonnegative
after a vector step and the last \(s\) rows of \(\lambda\) are
nonpositive (interlace downwards from the zero middle rows).

In the initial list
\[
 c_i\le B+\delta A_i\le2B.
\]

Thus \eqref{eq:kept-tail} gives
\(u_i\le A_i(1+\delta)+2\mathcal L B/k\) and
\[
 c_{i+1}\ge c_i-(\delta+2\mathcal L/k)B.
\]
Starting at
\((1-\epsilon)B\) or higher gives
\[
 c_i\ge(1-4\epsilon)B\ge A_i
\]
throughout that list.
Reinserting this comparison in \eqref{eq:kept-tail} gives
\[
 u_i\le A_i(1+\delta)+\mathcal L c_i/k,\qquad h_i=u_i-w_i\le\mathcal L c_i/k\le2\mathcal L B/k.
\]
In the tied list one has
\(c_i\le B+\delta A_i<(1-\Delta+\delta)A_i\).

Consequently \(u_i\le c_i+\mathcal L A_i/k\le A_i\).
The final degree satisfies
\[
 d\le c_s+\sum_i(C_i-n_i)\le c_s+(\mathcal L/k)\sum_i C_i.
\]
In fact, on the final upper rows the increments of the
restriction telescope by interlacing to at most \(c_s\)
through the last middle row. On the last \(s\) rows the
negative mass before restriction is at most \(\sum C_i\)
and afterward equals \(\sum n_i\). This proves the first
inequality even when some last rows vanish. Since
\(C_i\le 2a/P_i\), the cutoff for \(o_+\) holds.
The first rows of the reduced path at a tie and at the end
are at most \(c_i\) and \(c_s\) by interlacing, exactly the
cutoffs \(\lfloor K_i-m_0\rfloor\) and
\(\lfloor K_s-m_0\rfloor\). Finally
\[
 j_i=h_i\le\min(o_i,C_i)
\]

and \(j_+\le\min(o_+,b_+)\), giving legal levels. The map is
injective by the reconstruction above and preserves the
weight of each path.

\proofstep{Weight of discarded paths.}
Write \(\eta=1/k^2\) and \(\tau=\mathcal L/(4k)\). Use the following coordinates for one step,
\[
 \begin{aligned} x&=c_i+r, & y&=(A_i+r)(1+\delta),\\ v&=R_i+r, & z&=u_i-R_i+\delta(A_i+r).\end{aligned}
\]
Then
\(0\le v\le\min(x,y)\), \(z\ge0\), and
\[
 v(c_{i+1}+r)=v(x+y-v-z)
             =xy-(x-v)(y-v)-vz\le xy .
\]
If the first inequality of \eqref{eq:kept-tail} fails then
\(z>2\tau\max(x,y)\), since \(A_i\ge k^3\) and
\(\max(x,y)\le2\max(c_i,A_i)\).
For \(v\ge\min(x,y)/2\) the term \(vz\) exceeds \(\tau xy\).
For smaller \(v\) the other term is at least \(xy/4\ge\tau xy\).
Thus the step has the extra factor \(1-\tau\) on its upper
bound.

Also \(y\le e^{2\eta} A_i\) and
\(n_i+r\le e^{2\eta} C_i\) since \(\delta\le\eta\) and
\(r/C_i,r/A_i\le\eta\). A failure of the second inequality
gives
\[
 n_i+r\le (1-\tau)e^{2\eta} C_i,
\]
by
\(n_i<(1-4\tau)C_i\), \(r/C_i\le\eta\) and \(\eta\le\tau\).
Multiplying and cancelling the successive \(c_i+r\) gives
for any failed path, since \(\lambda_0\le c_s\) and
\(A+r\le e^\eta A\),
\[
 \frac{(\lambda_0+r)\prod_i(R_i+r)(n_i+r)}{A\prod_i A_iC_i}
 \le e^{(4s+1)\eta}(1-\tau).
\]
The \((r-1)\)-st power of the right side is at most
\(e^{2-\mathcal L/5}\). Indeed \(2s+1\le r\) and
\((r-1)/(4k)\ge1/5\).

\proofstep{Converting the path loss into a dimension loss.}
To turn this into a dimension bound, row \(s-i\) of
\(\lambda\) is nonnegative and at most \(R_i\).
Each remaining vector interlaces it with the next row,
each covector can only decrease a fixed row, and the
final restricted row is at most the corresponding full
row. Its lower bound follows by interlacing with the zero
middle rows, or with the nonnegative entries above them.
Rows with \(s<j<r-s\) vanish.

With \(V\) the standard representation, the desired occurrence is
\[
 V_\lambda\lhook\joinrel\longrightarrow \bigotimes_{\lambda_i>0}\operatorname{Sym}^{\lambda_i} V\,\otimes\!\bigotimes_{\lambda_i<0}\operatorname{Sym}^{-\lambda_i} V^*.
\]
Fill positive rows singly from the top, then negative
rows singly from the bottom, using Pieri at each step.
Consequently
\[
 d_\lambda\le D_r(\lambda_0)\prod_iD_r(R_i)D_r(n_i).
\]
Use \(z^{r-1}\le(r-1)!D_r(z)\le(z+r)^{r-1}\) and
\(D_r(z)\le D_k(z)\) on the unrestricted tensor. A single
failed path therefore has weight at most
\(e^{2-\mathcal L/5}T(A)\).

There are at most \((8a+1)^{(2s+3)(2s+1)}\) possible paths
for \(s\ge1\). At each of at most \(2s+3\) weights only the
top \(s+1\) and bottom \(s\) entries need be recorded.
Each has absolute value at most \(4a\), since at a vector or
covector step an individual absolute change is at most
the degree, and
\[
 A+\sum_i(A_i+C_i)\le4a
\]
by the
geometric progression of the \(P_i\), \(P_i\ge p\),
\(A\le a\) and \(A_i,C_i\le 2a/P_i\).

The same bound holds upon interlacing. Thus the log of the
number is at most \(30N_1^2\log a\), using \(s\le N_1\)
and \(a\ge25\). Increasing the absolute constant \(C_0\)
if necessary, this and the additive 2 in the exponent
are at most \(3\mathcal L/40\). The total failed weight
is at most \(e^{-\mathcal L/8}T(A)\). For \(s=0\)
there are no failures of \eqref{eq:kept-tail}.

\proofstep{The loss at the last label.}
For the last label put
\[
 J=\lfloor\mathcal L b_+/k\rfloor.
\]

The discarded monomials have last-coordinate exponent
greater than \(J\). If \(J<b_+\), subtract \(J+1\) from the last coordinate
of each discarded monomial. This bijects them with monomials
of degree \(b_+-J-1\), so their proportion is
\[
 \frac{D_k(b_+-J-1)}{D_k(b_+)}
 =\prod_{i=1}^r\left(1-\frac{J+1}{b_++i}\right)
 \le \exp[-r(J+1)/(b_++r)]\le e^{-\mathcal L/3},
\]
using \(b_+\ge k\) and \(k\ge 3\). For \(J\ge b_+\)
the discard is zero. Lastly
\(D_k(B-M)/D_k(B)\ge1-rM/B\) by the binomial product.
The two counts now give the lower bound in
\eqref{eq:ordinary-class} with coefficient
\begin{equation}\label{eq:rank-fraction}
 (1-e^{-\mathcal L/8})(1-e^{-\mathcal L/3})(1-rM/B).
\end{equation}
It exceeds \(1/2\) because \(\mathcal L\ge24\) and
\(kM/B\le1/30\). This completes the proof of
\eqref{eq:operational-seg} at the stated depths.

\end{proof}

\paragraph{Multiplication around the cycle.}

The ranks just estimated belong to individual segments. On closing the cycle the retained space has the product rank, so each discarded fraction must be charged before multiplying.

On the cycle let \(R_I\) be the rank of the retained digit
classes on segment \(I\), \(U\) the two-letter tensor space,
and \(g\) the number of neighboring edges. Denote by
\(\kappa_P\) the margin loss
\[
 \kappa_P=C_\kappa N\left(e^{-\mathcal L/8}
   +\frac{kQ(m_0+N_1 k)}{b'}+\frac{1}{Y}+\frac{kE_0}{a}\right),
\]
where \(C_\kappa\) is absolute.

\begin{lemma}[Product of retained ranks]\label{lem:cycle-retained-rank}
For each segment of the cycle impose the degree separation
and reserves of Lemma~\ref{hyp:retained-digit-degrees}.
The retained ranks satisfy
\[
 \prod_I R_I\ge(1-\kappa_P)\frac{\dim U}{g}.
\]

This is the digit-class loss to be added to the exceptional
fractions of Proposition~\ref{prop:segment-discard}.
\end{lemma}

\begin{proof}
\proofstep{The losses on neighboring segments.}

Let \(I\) index a segment, \(Q_I=p^{G_I}\), and let \(a_0\)
denote the actual degree at its starting even top. Summing
over the digits gives the full monomial dimensions at each
odd top, as in Section~\ref{sec:segment-window}. At the start,
\(B=\lfloor(a-E_0)/Q_I\rfloor\) and \(a\le a_0\le a+2a/(kY)\).
For every \(1\le j\le r\),
\[
 \frac{Q_I(B+j)}{a_0+j}\ge 1-z,\qquad
 z=\frac{2a/(kY)+E_0+Q}{a}.
\]
Indeed the difference between denominator and numerator
is at most \(2a/(kY)+E_0+Q_I\), and \(Q_I\le Q\).
In a nonvacuous lower bound \(z<1\), hence
\(Q_I^r D_k(B)/D_k(a_0)\ge1-rz\).
Every even start and every even end occur once.

Also
\(\prod_I Q_I=p^{n'}\) and the number of neighboring edges
is \(g=1+p^{n'}+\cdots+p^{rn'}\ge\prod_I Q_I^r\).
Multiplying the start factors and \eqref{eq:rank-fraction}
therefore gives
\begin{equation}\label{eq:retained-product}
 \prod_I R_I\ge(1-\kappa_P)\dim U/g,\qquad
 \kappa_P=C_\kappa N\left(e^{-\mathcal L/8}
   +\frac{kQ(m_0+N_1 k)}{b'}+\frac{1}{Y}+\frac{kE_0}{a}\right).
\end{equation}
Here \(C_\kappa\) is absolute. In detail, per segment use
\(rM/B\le2kQ(m_0+1)/a\), \(e^{-\mathcal L/3}\le
e^{-\mathcal L/8}\), and \(rz\le2/Y+k(E_0+Q)/a\).
For nonnegative losses whose sum is below one,
the product of their complements is at least one minus
their sum. Otherwise the assertion with enlarged
\(C_\kappa\) is automatic. This argument explains why a
fixed positive retained fraction per segment alone would
not give the required rank.

\proofstep{The specialization of the cycle parameters.}
For the cycle with parameters from \eqref{eq:retained-product}, take
\[
 E_0=C_{E1}\left(\frac{a}{p^2}+\frac{b'}{p}
       +\frac{\mathcal L N_1a}{pk}+\frac{\mathcal L b'}{k}\right)
       +C_{E2}k^2N_1Q^2c,\qquad
 m_0=C_m kN_1Qc+12c.
\]
In \eqref{eq:retained-product} take \(k^3Q^2cY\le b'\),
\(k^3QY\le b'\), \(kb'\le a\), \(\mathcal L Y b'\le a\),
\(8\log(NY)\le\mathcal L\le k\), and \(p\ge k^2Y^2\).
The exponential is at most \(1/Y\). The three parts of the
margin term, after division by \(N_1/Y\), are bounded by
\(C_m,12,1\) respectively.

After the same division each of
the four \(C_{E1}\) terms in \(kE_0/a\) is at most \(C_{E1}\).
The last term is at most \(C_{E2}\) since \(k^3Q^2cY\le b'\le a\).
It follows that
\(\kappa_P\le C_\kappa(C_m+4C_{E1}+C_{E2}+15)NN_1/Y\).
  All these estimates use the
same \(c,m_0\) and digit classes as the pairing tests.
\end{proof}
